
Reinforcement Learning from Human Feedback (RLHF) relies on preference datasets to align language models, yet whether these datasets encode exploitable geometric structure for alignment evaluation remains unexplored. This study tests whether aggregated human safety judgments induce clustering in semantic embedding space, enabling reusable benchmarks without per-model reward training. Through a three-hypothesis validation protocol applied to 160,800 chosen/rejected pairs from Anthropic's HH-RLHF dataset, we find: (1) the dataset contains genuine safety violations at 45.6\% base-rate (95\% CI: [41.3\%, 50.0\%], binomial p=0.0063); (2) human annotators achieve substantial inter-rater agreement (Cohen's κ=0.724, 95\% CI: [0.658, 0.791]) when applying explicit safety criteria; yet (3) RoBERTa-base embeddings exhibit no meaningful clustering (Cohen's d=0.034, p=0.797), with effect size 93\% below the target threshold despite statistical power exceeding 0.99. This negative result demonstrates a disconnect between human-detectable and embedding-detectable structure. Alignment violations span diverse semantic content (toxicity, misinformation, harmful instructions) that occupies overlapping embedding regions despite human-detectable differences. The systematic hypothesis decomposition localizes the failure at the representation level rather than dataset quality or annotation consistency, indicating that embedding-based alignment evaluation requires safety-specialized representations rather than general-purpose pretrained models.
